% What Survives a DELETE? - DASFAA 2027 Industrial Track (LNCS, single-blind, 12 pages)
% v0.7 - six models; industry-track review fixes. Only the artifact link (\artifacturl) remains to be filled.
% BEFORE SUBMISSION: put the artifact link in \artifacturl below (search for ARTIFACT-LINK).
\documentclass[runningheads]{llncs}
\usepackage[T1]{fontenc}
\usepackage{booktabs}
\usepackage{xcolor}
\usepackage{url}
\usepackage{array}
\usepackage{bbding}
\newcommand{\pending}[1]{\textcolor{red}{[PENDING: #1]}}
% ARTIFACT-LINK: replace the next line with  \newcommand{\artifacturl}{\url{https://github.com/USER/REPO}}
\newcommand{\artifacturl}{\url{https://github.com/dsyam22/erasure-residue-artifact}}

\begin{document}

\title{What Survives a DELETE? Measuring Erasure Residue in Change Data Capture and Temporal Tables}
\titlerunning{What Survives a DELETE?}

\author{Syam Sajit Dondapati\inst{1}\orcidID{0009-0004-0576-7152} \and
Varunbabu Suddapalli\inst{2}\orcidID{0009-0004-5384-6824}}
\authorrunning{S. S. Dondapati and V. Suddapalli}
\institute{Independent Researcher, NJ, USA\\
\email{dsyam22@gmail.com} \and
Independent Researcher, USA\\
\email{varunbabu.suddapalli@gmail.com}}

\maketitle

\begin{abstract}
The GDPR gives people a right to erasure, while audit rules push organisations
to record every change. Change data capture (CDC), system-versioned tables,
binary logs and logical replication serve the second need. On SQL Server 2022, PostgreSQL 17 and MariaDB 11.4 we erase
synthetic data subjects and measure what these features keep and what removing
it costs. Every feature kept the erased values in every trial. Retention
periods acted as lower bounds, cleanup runs that removed nothing were logged as
succeeded, temporal queries hid expired history rows that a direct read still
returned, and clearing a change stream ahead of a lagging consumer left every
erased subject in the downstream copy in 30 of 30 trials. We then asked six
language models to write the verification script. With the change-tracking
configuration withheld from the prompt, they certified an erasure that had not happened in 19--79\% of scripts per model;
given it, in at most 6\%. When querying the
catalog was optional, scripts that queried rarely failed (5 of 155) and scripts
that did not usually did (79 of 125). Scripts, raw results and a checklist are
released.

\keywords{Data erasure \and GDPR \and Change data capture \and Temporal tables
\and Logical replication \and Large language models}
\end{abstract}

% ---------------------------------------------------------------------------
\section{Introduction}

A database administrator who receives an erasure request under Article~17 of
the GDPR~\cite{gdpr} usually answers it with a \texttt{DELETE}, or with an
\texttt{UPDATE} that overwrites the personal fields so that foreign keys stay
intact. The row disappears, a query confirms it, and the request is closed.

The same database often serves auditors and downstream systems that need to
know what a record contained before it changed, through features that copy
rows as they change. Each is doing its job when it keeps a copy of a row that
has just been deleted.

We measure the gap between the two obligations. After a person's data has been
erased from the tables an application reads, where can a query still read it,
for how long, and what removes it? Nothing we report is a vendor bug, and much
of the behaviour is documented~\cite{mscdc,mstemporal,pgconfig}. What has been
missing is a measurement of its consequences for erasure. We make four
contributions.

\begin{enumerate}
\item \textbf{A baseline count of the residue} in change tables, history
tables, replication slots and binary logs on three engines
(Section~\ref{sec:rq1}), which fixes what there is to remove.
\item \textbf{What removes erasure residue, and at what cost}
(Section~\ref{sec:rq2}): retention acts as a lower bound, empty cleanup runs
report success, temporal retention hides rows hours before removing them, and
clearing a stream ahead of its consumer strands deletes, without an error in 20
of 30 trials.
\item \textbf{An evaluation of 2,249 LLM-written verification scripts} from
six models (Section~\ref{sec:rq3}), executed against databases whose residue is
known. Given the configuration, no model certified an unfinished erasure in
more than 6\% of scripts; without it, whether a model looked decided most
outcomes.
\item \textbf{A checklist and replication package} (Section~\ref{sec:checklist}).
\end{enumerate}

All data is synthetic, and the paper makes no claim about any production system
and no legal claim: keeping history for audit can be lawful. Our point is that
an operator asked to erase a person should know where copies remain.

% ---------------------------------------------------------------------------
\section{Background and Related Work}

\textbf{Change data capture.} SQL Server's CDC reads the transaction log and
writes one row per change into a change table that mirrors the source table's
columns: the full image for an insert or delete, before and after images for an
update~\cite{mscdc}. A cleanup job removes entries older than a retention
period, three days by default. A column added to the source later is not
captured, and a dropped column stays in the change table.

\textbf{System-versioned tables.} SQL Server moves the previous version of an
updated or deleted row into a history table, kept forever unless a retention
period is set; with one, a background task removes aged rows and temporal
queries filter them out~\cite{mstemporal}. MariaDB keeps history rows in the
table itself.

\begin{sloppypar}
\textbf{Logs and replication.} A PostgreSQL replication slot makes the server
retain write-ahead log until a consumer confirms it, by default without
limit~\cite{pgconfig}; whether a delete record carries the whole old row
depends on \texttt{REPLICA IDENTITY}~\cite{pgaltertable}. MariaDB's row-based
binary log records row images for replication.
\end{sloppypar}

\textbf{Related work.} GDPRbench~\cite{shastri2020} measures the cost of
retrofitting GDPR features such as timed deletion into Redis and PostgreSQL,
and GDPRizer~\cite{agarwal2022} finds all rows belonging to one person across a
legacy schema; both work on live tables, not on copies kept by change tracking.
DELF~\cite{delf2020} found thousands of omissions in
application-level deletion logic; we look one layer down, at copies the engine
makes without application code. Chakraborty et
al.~\cite{chakraborty2025} formalise erasure when erased values can be inferred
from dependencies; our copies are explicit. Lethe~\cite{sarkar2020} bounds the
time from a logical delete to physical removal in LSM-trees; we find that
relational retention settings bound it only from below. K9db~\cite{k9db} and
query-language support for timely deletion~\cite{sarkar2022} change the system
to make deletion correct; we measure unmodified engines. Other features keep
history by design: SQL Server ledger tables block deletes or keep every prior
version~\cite{msledger}, and Oracle Flashback Time Travel archives row changes
for a retention period~\cite{oraclefta}; Delta Lake offers time travel over
earlier table versions and deletes their files only after a retention
threshold~\cite{delta}. Crypto-shredding, which encrypts each
subject's data under its own key and erases it by deleting the key, is a common
answer for append-only logs~\cite{verraes2019}.

Pipelines such as Debezium emit each delete with the row's prior state and a
tombstone for Kafka compaction~\cite{debezium}, so residue recurs in the topic;
we cover only the step before it. Backup
purging~\cite{lenard2024} and database carving~\cite{wagner2015} concern copies
a \texttt{SELECT} cannot reach; ours it can. LLM work on
schema mapping~\cite{fu2024} and compliance Text-to-SQL~\cite{atf2025} does not
run model-written SQL to find residual personal data.

% ---------------------------------------------------------------------------
\section{Study Design}
\label{sec:design}

\begin{description}
\item[RQ1 (residue).] After an erasure, where are a data subject's original
values still readable, and in how many copies?
\item[RQ2 (remediation).] What removes each residue, how long does it take, and
what does it cost elsewhere?
\item[RQ3 (verification).] Can an LLM write a script that finds every surface
holding residue?
\end{description}

\textbf{Environment and data.} We use SQL Server 2022 Developer Edition (build
16.0.4295.3, CU27), PostgreSQL 17.11 and MariaDB 11.4.13 in Docker containers on
one Apple Silicon laptop; the first two run as x86-64 images under emulation.
Each trial starts from a new database. Tables hold synthetic customers with a
name, an e-mail address and a telephone number, at $10^3$, $10^5$ and $10^6$
rows ($10^3$ and $10^5$ on MariaDB). In each trial 100 subjects are erased by a
\texttt{DELETE} or by an \texttt{UPDATE} that overwrites the personal fields.

\textbf{Residue.} A subject has residue on a surface when a \texttt{SELECT}, or
the engine's own log reader for a slot or binary log, returns one of the
subject's original values after the erasure has committed and the capture
process has recorded it. We poll until it has, instead of sleeping for a fixed
time. We do not inspect storage pages. RQ1 outcomes are deterministic; we repeated
each configuration ten times (three on MariaDB) to rule out timing effects and
report them as identical, not as proportions. Timings are ranges over trials.

\textbf{Scenarios.} RQ1 covers SQL Server CDC and temporal residue (3 scales
$\times$ 2 methods, 60 trials), PostgreSQL slots under both replica-identity
settings (120), MariaDB history and binary log (12), and schema changes under
CDC (10). RQ2 covers CDC cleanup timing at three retention settings on an idle
database and on one with an unrelated change every 30 seconds (60), temporal
retention over a day (one run, then five staggered runs), downstream copies on
SQL Server (20) and PostgreSQL (10), per-subject removal and read access (3,
plus the MariaDB trials), and four restore sequences (40).

\textbf{LLM task.} We built nine SQL Server databases (cases i1--i9), each with
one erased subject: seven hold residue on a known set of surfaces (change table, history
table, both, a dropped column's old values, or a downstream copy) and two are
clean. A model receives the subject's original values, a description of the
erasure and (1, \emph{withheld}) the table definitions only; (2) these and a listing of the
change-tracking configuration; (3) both and a hint (``consider change-tracking
features (CDC change tables, temporal history tables) and downstream copies of
the data''); or (4, \emph{discovery}) what condition~1 gives plus one read-only
query of its own, whose output it sees before writing the script; or (5,
\emph{optional discovery}) the same with up to three queries that the model may
run or skip; or (6, \emph{agentic}) one to three queries with errors shown and
one repair of a failing script. No discovery instruction names a
change-tracking feature. The model returns a
read-only T-SQL script that reports, per object checked, how many rows still
hold the data. Scripts run under a read-only login. Ground truth is the set of
surfaces each case was built to hold; an exhaustive search run at the moment
each script runs confirms it, and no model judges another. Case i5 (one-day temporal retention) holds residue only
between expiry and removal, so its scripts were run inside that window. We report
\emph{false clean} (data reported erased where residue exists), execution
without error and false alarms on the clean cases; the share of surfaces found
per script is in the artifact. Output not in the requested form counts as finding nothing, not as false
clean; a script that errs but returns a result is scored on that result. Two
baselines bracket the models: a \emph{naive} check of the source table plus a
temporal query, and the exhaustive search, a schema-agnostic scan of every table
that has a column of the searched name (source, change, history and copy
tables); we release it as a reference verifier, valid where column names are
kept in copies. Models: Gemma 4 E4B (local), DeepSeek Flash and V4 Pro
(reasoning off), and Gemini 3.1 Pro, GPT-6.1 Sol and Claude Haiku 5.5 at their
lowest reasoning settings, queried in October 2026, at temperature 0.7 to
expose variance (GPT and Haiku: provider default, which their APIs fix). Ten
trials per model, case and condition; five for GPT and Gemini (budget, Gemini's
daily quota); condition 3 for the first three models only, 5--6 for hosted
models; condition 6 with five trials (GPT, Gemini: three; DeepSeek: no i5); V4
Pro with reasoning on for conditions 1 and 4, five trials.

% ---------------------------------------------------------------------------
\section{RQ1: Where Erased Data Persists}\label{sec:rq1}

\textbf{SQL Server.} Table~\ref{tab:residue} gives the result, the same in all
60 trials. The source tables held none of the 100 subjects' values; the change
table held two copies of each and the history table one, which a temporal query
with \texttt{FOR SYSTEM\_TIME ALL} returned. The erasure adds a copy: before the
\texttt{DELETE} the change table held one image of the row, and the delete
record is a second. Overwriting does not help either, since the update's
before-image and the closed history version carry the original values.

\begin{table}[t]
\caption{Readable copies of an erased subject's values on SQL Server, identical
in all 60 trials. The change-table copies are the insert image and the record
of the erasure (delete record or the update's before image).}\label{tab:residue}
\centering
\small
\begin{tabular}{lccl}
\toprule
Surface & \texttt{DELETE} & Overwrite & What the copy is \\
\midrule
Source table & 0 & 0 & -- \\
CDC change table & 2 & 2 & insert image, erasure record \\
Temporal history table & 1 & 1 & closed version \\
\midrule
Subjects with residue & 100/100 & 100/100 & \\
\bottomrule
\end{tabular}
\end{table}

\textbf{PostgreSQL.} What the erasure puts into the change stream depends on
replica identity. With \texttt{FULL}, often set for tables without a suitable
key, each of the 100 delete or update records carried the complete old row;
with the default identity, only the key. This is the exposure of a working
pipeline: the record waits in the slot until the consumer confirms it, and is
then delivered downstream. Our slot, created before the load and never read
(the extreme case of an abandoned consumer), also retained the original inserts,
so all 100 subjects were readable through it in all 120 trials; once a consumer
had read it, nothing was. At $10^6$ rows the unread slot held
back about 187\,MB of log, with no limit configured by default~\cite{pgconfig}.

\textbf{MariaDB.} With system versioning and the row-based binary log enabled,
all 100 erased subjects remained in the table's history in all 12 trials, and
the binary log held three copies of each subject's original values, for both
methods and both scales: the insert, then a delete logged as an update whose before and
after images both carry the original values (the row is closed, not removed),
or an overwrite's before image plus the history row it writes.

\textbf{Schema changes.} In all ten trials a column added under CDC was not
captured, and a dropped column stayed in the change table with its values
readable~\cite{mscdc}: a change log can record too little for an audit and too
much for an erasure.

% ---------------------------------------------------------------------------
\section{RQ2: Removing the Residue}\label{sec:rq2}

\subsection{CDC Cleanup: The Retention Period Is a Lower Bound}

We deleted one subject and started the cleanup job every 30 seconds
(Table~\ref{tab:cleanup}). On a database receiving unrelated changes, the
subject's rows went within about a minute of the retention period. On an idle
database, removal moved in five-minute steps, and the delete record outlived the
insert image by a further five minutes in all 30 idle trials. The documentation
explains the steps: cleanup sets its low-water mark from the newest entry in the
LSN-to-time mapping, not from the clock~\cite{mscdc}, and on an idle database
that entry can be minutes old.

\begin{table}[t]
\caption{Minutes from erasure until the subject's rows left the CDC change
table (range over 10 trials; polled every 30\,s).}\label{tab:cleanup}
\centering
\small
\begin{tabular}{llcc}
\toprule
Retention & Database & Insert image gone & Delete record gone \\
\midrule
1 min & active & 1.6--2.1 & 1.6--2.1 \\
1 min & idle & 5.6--5.7 & 10.6--11.1 \\
5 min & active & 5.6--6.2 & 5.6--6.2 \\
5 min & idle & 5.6--5.7 & 10.6--11.1 \\
15 min & active & 15.6--16.2 & 15.6--16.2 \\
15 min & idle & 15.6--16.2 & 20.6--21.2 \\
\bottomrule
\end{tabular}
\end{table}

Two details matter more than the overshoot. The last copy to go is the delete record, the one the erasure created. And
a run that removes nothing looks like one that does: counting runs started after
the period had passed that left the insert image in place, excluding the run
that removed it, there were eight in every trial at the one-minute setting on an
idle database, and the job history records each as succeeded.

\subsection{Temporal Retention: Hidden Before It Is Gone}

We set a one-day retention period, deleted 100 subjects and read the history
table directly and through a temporal query. A first run, read every 15
minutes, hid all 100 from the temporal query at the first reading after expiry
and removed them 32 to 48 minutes after expiry. In five more databases, erased
seven minutes apart and read every minute, the temporal query returned all 100
subjects at the last reading before expiry and none at the first after it. A direct read of the history table still
returned all 100 for 3.9 to 4.4 hours after expiry, until the subjects of all
five databases disappeared between the same two readings, consistent with one
background pass rather than removal at each run's expiry; the documentation
gives no bound~\cite{mstemporal}. Meanwhile a check through a temporal query reports
the data erased while a \texttt{SELECT} on the history table returns it.

\subsection{Downstream Copies: Clearing One Surface Breaks Another}

A downstream table was fed by a consumer that applies deletes; when expected
changes are gone, one design moves its read window forward, the other lets the
engine raise an error. While the consumer kept up, deletes reached the copy in every trial.
We then deleted ten subjects with the consumer paused, shortened retention so
that cleanup removed their change records (1.6 to 2.2 minutes), and ran the
consumer. In all 20 SQL Server trials (ten per design) the ten subjects
remained in the copy.
The first design reported nothing; the second got error 313, which signals the
break but leaves the copy equally stale. On PostgreSQL, with the subscription paused,
we deleted ten subjects, dropped and recreated the slot (a step operators take
when an idle slot holds back log) and resumed. In 10 of 10 trials the subjects
remained in the copy, later inserts arrived normally, and the subscription
reported no errors. The step that clears the change stream destroys the only
signal telling a downstream copy to delete.

\subsection{Removing One Person, and Who Can Read the Residue}

We tested per-subject removal on SQL Server in three trials, each on a
1,000-row table from which ten subjects had been deleted (1,010 change rows),
and on MariaDB in its 12 trials; all agreed. As database owner, a \texttt{DELETE} on the SQL
Server change table removed the subject's two entries and left other entries
and the change functions intact, but a second capture instance on the same
table kept its own two copies. The change table then held exactly 1,008 of
its 1,010 rows; we did not test the change functions over windows that span the
deleted rows.
This works but is outside documented use: the
documentation recommends reading change tables only through the CDC
functions~\cite{mscdcct} and offers no per-row removal. A \texttt{DELETE} on the history table was
refused while versioning was on (error 13560); with versioning off it
succeeded, versioning went back on with a consistency check, and other
subjects' history was untouched.

\begin{sloppypar}
MariaDB rejected both per-row forms we tried, \texttt{DELETE ... FOR
SYSTEM\_TIME ALL WHERE} and \texttt{DELETE HISTORY ... WHERE}, as syntax errors,
also with \texttt{system\_versioning\_\allowbreak alter\_\allowbreak history = KEEP}. Its
\texttt{DELETE HISTORY ... BEFORE SYSTEM\_TIME} removed the erased subjects'
history together with every other row's. We rotated the binary log before each
purge, so the erased values were never in the active file; even so, in 7 of 12
trials the first \texttt{PURGE BINARY LOGS} left the file in place, reporting
only a note (1375) that it ``may be needed for crash recovery''; a second call
removed it. The SQL Server remedy for a dropped column, a new capture instance
with the old one disabled, discarded every entry of the old change table
(1,008: 1,000 inserts and ten deletes, less the two removed above).
\end{sloppypar}


Access differs by surface. On SQL Server a member of \texttt{db\_datareader}
read the change table and the history table. A user granted \texttt{SELECT} on
the source tables only was refused on both, but read the erased subjects'
entries through the CDC query function, because the capture instance had no
gating role. On PostgreSQL a role with \texttt{pg\_read\_all\_data} was refused
the slot; reading it needs the \texttt{REPLICATION} attribute.

\subsection{After a Restore}

That restoring an older backup returns erased rows is known~\cite{lenard2024};
we looked at change capture after four restore sequences (ten trials each, all
agreeing). \texttt{KEEP\_CDC} kept the change table with the erased subject's
row but restored neither the capture nor the cleanup job, so that residue has no
time bound until someone adds them. A default restore disabled CDC. Temporal
retention cleanup was off after a direct full restore and on when a log restore
came last.

% ---------------------------------------------------------------------------
\section{RQ3: LLM-Written Verification}\label{sec:rq3}

\begin{table}[t]
\caption{LLM-written verification scripts (2,249). Top: \emph{false clean}
(residue-case scripts reporting the data erased; 7 of the 9 cases are residue
cases). Bottom: scripts that did not run. Only Gemma reported residue for a
clean case (once each in conditions 2 and 3).}\label{tab:llm}
\centering
\footnotesize
\setlength{\tabcolsep}{4pt}
\begin{tabular}{@{}lcccccc@{}}
\toprule
Model & 1 Withheld & 2 Given & 3 Hint & 4 Disc. & 5 Optional & 6 Agentic \\
\midrule
\multicolumn{7}{@{}l}{\emph{False clean}} \\
Gemma 4 E4B & 50/70 & 0/70 & 1/70 & 21/70 & -- & -- \\
DeepSeek Flash & 50/70 & 1/70 & 0/70 & 5/70 & 46/70 & 4/30 \\
DeepSeek V4 Pro & 55/70 & 0/70 & 0/70 & 16/70 & 31/70 & 3/30 \\
V4 Pro, reasoning & 23/35 & -- & -- & 3/35 & -- & -- \\
Claude Haiku 5.5 & 13/70 & 4/70 & -- & 3/70 & 2/70 & 0/35 \\
GPT-6.1 Sol & 8/35 & 0/35 & -- & 2/35 & 0/35 & 1/21 \\
Gemini 3.1 Pro & 18/35 & 0/35 & -- & 3/35 & 5/35 & 5/21 \\
Naive check & 4/7 & & & & & \\
Exhaustive search & 0/7 & & & & & \\
\midrule
\multicolumn{7}{@{}l}{\emph{Did not run}} \\
Gemma 4 E4B & 15/90 & 14/90 & 18/90 & 38/90 & -- & -- \\
DeepSeek Flash & 6/90 & 10/90 & 32/90 & 5/90 & 13/90 & 2/40 \\
DeepSeek V4 Pro & 5/90 & 8/90 & 19/90 & 4/90 & 10/90 & 3/40 \\
V4 Pro, reasoning & 1/45 & -- & -- & 3/45 & -- & -- \\
Claude Haiku 5.5 & 7/90 & 5/90 & -- & 4/90 & 3/90 & 0/45 \\
GPT-6.1 Sol & 0/45 & 0/45 & -- & 0/45 & 1/45 & 0/27 \\
Gemini 3.1 Pro & 2/45 & 0/45 & -- & 0/45 & 0/45 & 0/27 \\
\bottomrule
\end{tabular}
\end{table}

\textbf{With the configuration withheld, the models split} (Table~\ref{tab:llm}).
The schema listing omits the change and history tables, so condition~1 measures
what the prompt withholds as much as the model; its rates are not general
failure rates. Gemma and the two DeepSeek models
reported the data erased in 155 of 210 residue-case scripts (71--79\% per
model), as did the naive check in 4 of 7 cases. Haiku, GPT and Gemini did so in
19\%, 23\% and 51\%: many of their scripts (69 of 70, 27 of 35 and 15 of 35)
queried the catalog, CDC or history tables instead of only the listed tables. Failures
cluster by case, so we treat model differences as descriptive: with the case
as the unit (exact sign-flip test over the seven residue cases, smallest
possible $p = 0.016$), Haiku and GPT differ from each older model at
$p = 0.06$--$0.09$.

\textbf{With the configuration, the wrong answer almost disappeared}: 5 of 350
scripts across the six models (4 from Haiku). The hint did less and, against
condition~2, lowered execution, for Flash from 80 to 58 of 90 scripts.

\textbf{Made to look first, most models found most of the residue}: Flash,
V4 Pro and Gemma reported the data erased in 5, 16 and 21 of 70 scripts, the
newer models in 2--3. For the older models the cause was not looking: in 40 of
their 42 false-clean scripts the case had a change table that the model's own
query neither mentioned nor returned.

\textbf{When looking was optional, looking decided the outcome.} With up to
three optional queries, Flash and V4 Pro queried in 4 and 23 of their 70
residue-case scripts and reported the data erased in 46 and 31; Haiku, GPT and
Gemini queried in 65 of 70, 35 of 35 and 28 of 35, and did so in 2, 0 and 5
(case-level $p = 0.03$ for Haiku and GPT against each DeepSeek model). Across
the five models, 5 of the 155 scripts that queried reported the data erased,
against 79 of the 125 that did not.

\textbf{CDC-only cases were the newer models' blind spot.} Of their 64
false-clean scripts over all conditions, 56 were for i1 and i6, whose only
residue is in a CDC change table. \textbf{Reasoning did not clearly help}: V4
Pro with reasoning on certified an unfinished erasure in 23 of 35 scripts with
the configuration withheld and 3 of 35 in discovery, against 16 of 70 without reasoning
(Fisher's exact test, $p = 0.11$; underpowered). \textbf{An agent loop did not
remove the failures}: with one to three required queries, errors shown and one
repair, Flash and V4 Pro failed in 4 and 3 of 30 scripts (1 and 7 on the same
cases and trials with one query), Haiku in 0 of 35, GPT in 1 of 21 and Gemini
in 5 of 21.
\textbf{Execution} separated the generations: 2--42\% of the older models'
scripts in a cell failed to run, 0--8\% of the newer models'.

The first author labelled 72 randomly drawn scripts of the first three models
(five per case from conditions 1--3, three from 4--5) without seeing the
automated verdict; the labels agreed with it for 43 of 45 and 27 of 27 (Cohen's
$\kappa = 0.93$ and $1.0$). The hosted-model runs cost about US\$11 in total
(DeepSeek \$3.60, Google \$3.68, OpenAI \$2.66, Anthropic \$0.66, from the
providers' billing).

% ---------------------------------------------------------------------------
\section{Practitioner Checklist and Discussion}\label{sec:checklist}

Table~\ref{tab:checklist} turns the measurements into a procedure: each surface
needs its own check and remedy.

\begin{table}[t]
\caption{Where to look after an erasure, how to clear it, and what clearing it
costs. Entries follow our measurements, except the remedies for downstream
copies and restored databases, which follow the documentation.}\label{tab:checklist}
\centering
\footnotesize
\setlength{\tabcolsep}{4pt}
\begin{tabular}{@{}>{\raggedright}p{1.75cm}>{\raggedright}p{3.1cm}>{\raggedright}p{3.1cm}>{\raggedright\arraybackslash}p{3.1cm}@{}}
\toprule
Surface & Check & Clear & Side effect \\
\midrule
CDC change table & Query each capture instance's change table by key, not job status & Delete the subject's rows in every instance (undocumented; checked by row counts only), or shorten retention and run cleanup & Retention clears everyone's history; a lagging consumer never sees the delete \\ \addlinespace[2pt]
Temporal history & Query the history table directly, not through a temporal query & Versioning off, delete the subject's rows, versioning on & No history recorded while off; retention hides rows before removing them \\ \addlinespace[2pt]
Dropped column & List each instance's captured columns & New capture instance, disable the old & Old instance's history for every row is dropped \\ \addlinespace[2pt]
Replication slot & List slots and the log they hold; check replica identity & Let the consumer read, or drop the slot & A dropped slot discards deletes not yet applied \\ \addlinespace[2pt]
MariaDB history, binary log & \texttt{FOR SYSTEM\_TIME ALL}; decode the binary log & \texttt{DELETE HISTORY}; purge logs, then list them & Removes every row's history; a purge may leave the file \\ \addlinespace[2pt]
Downstream copy & Compare keys with the source & Reconcile against the source & Must be scheduled \\ \addlinespace[2pt]
Restored database & Check capture and cleanup jobs and temporal retention & Re-add jobs; re-enable retention & -- \\
\bottomrule
\end{tabular}
\end{table}

Two further lessons concern order and verification. Clear downstream copies, or
confirm that their consumers have caught up, before clearing the stream that
feeds them. Give a model the change-tracking configuration, not just the
schema, or make it query the catalog before writing; check change tables
explicitly, where even the newer models failed; and treat a script that reports
nothing as unverified until it covers
every surface in Table~\ref{tab:checklist}; the exhaustive search is a sound
reference where copies keep column names (renamed, pseudonymised or packed
columns escape it). Restrict the CDC query functions with a gating role, since
\texttt{SELECT} on the source table is otherwise enough to read erased rows
through them. Prevention is cheaper than removal: leave personal columns out of a capture
instance's column list, pseudonymise them before they reach a captured table,
or encrypt them per subject so that deleting a key erases every copy at
once~\cite{verraes2019}. We did not measure these. For engine designers, the
measurements argue for a supported way
to remove one person's entries from every change and history structure at once,
and for a consumer-visible marker when changes are removed before being read.

% ---------------------------------------------------------------------------
\section{Threats to Validity}

\textbf{Environment.} SQL Server and PostgreSQL ran under emulation on one
laptop, so absolute times are not production figures; we rely on the ordering
and step structure in Table~\ref{tab:cleanup}. The temporal removal delay
depends on when a background task runs and is the result most likely to differ
on native hardware or a managed service. \textbf{Generality.} We tested one version of
each engine with synthetic data, not managed services such as Azure SQL
Database, where capture and cleanup are scheduled differently, and not an
external pipeline such as Debezium with Kafka. \textbf{LLM study.} Nine cases cannot represent real schemas. The cases
are new databases, though CDC verification queries are common in public
material. Scripts came from one prompt wording per condition. Hosted models change behind a fixed name, so our
results describe the versions served in October 2026; GPT and Haiku ran at their
providers' default temperature, and GPT and Gemini with five trials, not ten.
The hand labels cover the first three models only.

% ---------------------------------------------------------------------------
\section{Conclusion}

Change data capture, system-versioned tables, binary logs and logical
replication keep copies of erased rows by design, and they did so in every
trial on all three engines, for both ways of erasing. Each surface needs a different removal; only SQL Server permitted
per-subject removal, outside documented use or with versioning off; and removing
copies in the wrong order leaves a downstream copy stale, usually without an
error. All six models wrote a correct check when told which change-tracking
features are in use. Without that, what mattered was whether a model looked;
the newer models mostly did, but still missed CDC change tables. A verification
workflow should supply the configuration or require the look, and check change
tables explicitly.

\subsubsection*{Data and Artifact Availability.} Scripts, raw results, prompts
and every model-written script are available at \artifacturl.

\subsubsection*{Use of Generative AI.} Large language models are the object of
study in Section~\ref{sec:rq3}. In addition, the authors used an AI assistant
(Claude, Anthropic) to help write the experiment scripts and to draft and edit
the text; Claude Haiku 5.5, an evaluated model, comes from the same provider
and was scored like the others. The authors ran all experiments, checked every
number against the raw results, and are responsible for the content.

\subsubsection*{Disclosure of Interests.} The authors have no competing interests
to declare that are relevant to the content of this article.

\begin{thebibliography}{99}
\bibitem{gdpr} Regulation (EU) 2016/679 of the European Parliament and of the
Council (General Data Protection Regulation), Article 17. Official Journal of
the European Union L 119 (2016)
\bibitem{shastri2020} Shastri, S., Banakar, V., Wasserman, M., Kumar, A.,
Chidambaram, V.: Understanding and benchmarking the impact of GDPR on database
systems. Proc. VLDB Endow. \textbf{13}(7), 1064--1077 (2020)
\bibitem{chakraborty2025} Chakraborty, V., Kaminsky, Y., Mehrotra, S., Naumann,
F., Nawab, F., Pappachan, P., Sadoghi, M., Venkatasubramanian, N.: Meaningful
data erasure in the presence of dependencies. Proc. VLDB Endow.
\textbf{18}(10), 3435--3448 (2025)
\bibitem{sarkar2020} Sarkar, S., Papon, T.I., Staratzis, D., Athanassoulis, M.:
Lethe: A tunable delete-aware LSM engine. In: Proc. ACM SIGMOD International
Conference on Management of Data, pp. 893--908 (2020)
\bibitem{agarwal2022} Agarwal, A., George, M., Jeyaraj, A.R., Schwarzkopf, M.:
Retrofitting GDPR compliance onto legacy databases. Proc. VLDB Endow.
\textbf{15}(4), 958--970 (2022)
\bibitem{lenard2024} Lenard, B., Rasin, A., Scope, N., Al Johani, T.: The
intersection of compliance, databases, and IT operations. In: Proc. 36th
International Conference on Scientific and Statistical Database Management
(SSDBM) (2024)
\bibitem{wagner2015} Wagner, J., Rasin, A., Grier, J.: Database forensic
analysis through internal structure carving. Digital Investigation
\textbf{14}, S106--S115 (2015)
\bibitem{fu2024} Fu, S.D., Chen, X.: Compound schema registry.
arXiv:2406.11227 (2024)
\bibitem{atf2025} Atf, Z., Lewis, P.R.: ScenarioBench: Trace-grounded
compliance evaluation for Text-to-SQL and RAG. arXiv:2509.24212 (2025)
\bibitem{mscdc} Microsoft: About change data capture (SQL Server).
\url{https://learn.microsoft.com/en-us/sql/relational-databases/track-changes/about-change-data-capture-sql-server},
accessed October 2026
\bibitem{mstemporal} Microsoft: Manage retention of historical data in
system-versioned temporal tables.
\url{https://learn.microsoft.com/en-us/sql/relational-databases/tables/temporal/manage-retention},
accessed October 2026
\bibitem{mscdcct} Microsoft: cdc.<capture\_instance>\_CT (Transact-SQL).
\url{https://learn.microsoft.com/en-us/sql/relational-databases/system-tables/cdc-capture-instance-ct-transact-sql},
accessed October 2026
\bibitem{msledger} Microsoft: Ledger overview (SQL Server 2022).
\url{https://learn.microsoft.com/en-us/sql/relational-databases/security/ledger/ledger-overview},
accessed October 2026
\bibitem{oraclefta} Oracle: Track table changes with Flashback Time Travel.
\url{https://docs.oracle.com/en-us/iaas/autonomous-database-serverless/doc/flashback-time-travel-autonomous.html},
accessed October 2026
\bibitem{verraes2019} Verraes, M.: Eventsourcing patterns: Crypto-shredding.
\url{https://verraes.net/2019/05/eventsourcing-patterns-throw-away-the-key/}
(2019), accessed October 2026
\bibitem{pgconfig} PostgreSQL Global Development Group: PostgreSQL 17
documentation, Replication (max\_slot\_wal\_keep\_size).
\url{https://www.postgresql.org/docs/17/runtime-config-replication.html},
accessed October 2026
\bibitem{pgaltertable} PostgreSQL Global Development Group: PostgreSQL 17
documentation, ALTER TABLE (REPLICA IDENTITY).
\url{https://www.postgresql.org/docs/17/sql-altertable.html}, accessed October
2026
\bibitem{delf2020} Cohn-Gordon, K., Damaskinos, G., Neto, D., Cordova, J.,
Reitz, B., Strahs, B., Obenshain, D., Pearce, P., Papagiannis, I.: DELF:
Safeguarding deletion correctness in online social networks. In: 29th USENIX
Security Symposium (2020)
\bibitem{k9db} Albab, K.D., Sharma, I., Adam, J., Kilimnik, B., Jeyaraj, A.,
Paul, R., Agvanian, A., Spiegelberg, L., Schwarzkopf, M.: K9db:
Privacy-compliant storage for web applications by construction. In: 17th USENIX
Symposium on Operating Systems Design and Implementation (OSDI), pp. 99--116
(2023)
\bibitem{sarkar2022} Sarkar, S., Athanassoulis, M.: Query language support for
timely data deletion. In: Proc. 25th International Conference on Extending
Database Technology (EDBT) (2022)
\bibitem{delta} Armbrust, M., Das, T., Paranjpye, S., et al.: Delta Lake:
High-performance ACID table storage over cloud object stores. Proc. VLDB Endow.
\textbf{13}(12), 3411--3424 (2020)
\bibitem{debezium} Debezium Community: Debezium connector for PostgreSQL, Data
change events.
\url{https://debezium.io/documentation/reference/stable/connectors/postgresql.html},
accessed October 2026
\end{thebibliography}
\end{document}
